\section{模型评价、适用边界与发布结论}

\subsection{方法层面的评价}

密度贪心法的优点是实现简单、扫描过程可解释，排序后求解成本低；其局限是局部密度顺序不能保证
$0$--$1$ 组合的全局最优。穷举法在有限搜索空间内给出精确参照，适合本演示的已知小规模实例，
但复杂度为 $O(2^n)$，不能据此推荐为大规模求解器。

两个完整实例呈现了互补现象：在 \texttt{known-gap} 中，贪心选择受到不可分割性影响而低于穷举值；
在 \texttt{known-tie} 中，两法得到相同结果。这说明单一平均值会掩盖实例差异，因此正文同时保留原始实例图和
Metric 表。但样本过少且包含故意失败，本文不能比较稳定性、运行时间分布或一般最优性概率。

\subsection{发布边界}

本演示的贡献不在算法创新，而在可核验的发布链：
\begin{enumerate}
  \item 正式选择与 evidence freeze 共同绑定精确的运行清单哈希；
  \item 重复数字由程序生成宏和表格行，摘要、正文与表格共享同一来源；
  \item 正式图只读取已保存观测和 figure CSV，显示修改不触碰 Metric；
  \item 单位、样本、聚合、数值或正式引用一旦改变，发布流程停止并要求科学修订；
  \item 编译成功后仍需检查字体、文本层、匿名、关键数字并逐页查看栅格图，才可进入人工审阅。
\end{enumerate}

因此，本文只得出一个开发结论：冻结输入、确定性生成、脚本支持图表和实际 PDF 复核可以组成小而完整的
数学建模发布闭环。它不意味着 D4 合成算法已经通过真实竞赛验证，也不意味着演示中的人工采纳夹具构成批准。
